\documentclass[conference]{IEEEtran}
\IEEEoverridecommandlockouts

% Configuración para español
\usepackage[spanish,es-nodecimaldot,es-noshorthands]{babel}
\usepackage[utf8]{inputenc}
\usepackage[T1]{fontenc}

% Paquetes estándar IEEE
\usepackage{cite}
\usepackage{amsmath,amssymb,amsfonts}
\usepackage{algorithmic}
\usepackage{graphicx}
\usepackage{textcomp}
\usepackage{xcolor}
\usepackage{float}
\usepackage{listings} % Para código

% ---------- Configuración de listados de código ----------
\definecolor{codegray}{rgb}{0.45,0.45,0.45}
\definecolor{codeblue}{rgb}{0.10,0.10,0.60}
\lstset{
	basicstyle=\ttfamily\scriptsize,
	keywordstyle=\color{codeblue}\bfseries,
	commentstyle=\color{codegray}\itshape,
	stringstyle=\color{codegray},
	numbers=left,
	numberstyle=\tiny\color{codegray},
	numbersep=5pt,
	frame=single,
	breaklines=true,
	breakatwhitespace=true,
	showstringspaces=false,
	tabsize=2,
	captionpos=b,
	literate=%
	{á}{{\'a}}1 {é}{{\'e}}1 {í}{{\'i}}1 {ó}{{\'o}}1 {ú}{{\'u}}1
	{Á}{{\'A}}1 {É}{{\'E}}1 {Í}{{\'I}}1 {Ó}{{\'O}}1 {Ú}{{\'U}}1
	{ñ}{{\~n}}1 {Ñ}{{\~N}}1 {¿}{{?`}}1 {¡}{{!`}}1
}

% ---------- Espacio reservado para imágenes faltantes ----------
\newcommand{\imgph}[2]{%
	\fbox{\begin{minipage}[c][#2][c]{0.92\columnwidth}
			\centering\textcolor{codegray}{\footnotesize [Espacio para imagen: #1]}
	\end{minipage}}%
}

\def\BibTeX{{\rm B\kern-.05em{\sc i\kern-.025em b}\kern-.08em
	T\kern-.1667em\lower.7ex\hbox{E}\kern-.125emX}}

\begin{document}

\title{Práctica 4: Proceso de Arranque (Boot)}

\author{
	\IEEEauthorblockN{Belmont Muñoz Samuel}
	\IEEEauthorblockA{\textit{Ingeniería en Computación} \\
		\textit{Facultad de Ingeniería}\\
		No. Cuenta: 317318289}
	\and
	\IEEEauthorblockN{Domínguez Palacios Jesús Alejandro}
	\IEEEauthorblockA{\textit{Ingeniería en Computación} \\
		\textit{Facultad de Ingeniería}\\
		No. Cuenta: 421066205}
	\and
	\IEEEauthorblockN{Dorantes Mejía Emmanuel}
	\IEEEauthorblockA{\textit{Ingeniería en Computación} \\
		\textit{Facultad de Ingeniería}\\
		No. Cuenta: 313004766}
	\and
	\IEEEauthorblockN{García Olvera Emiliano}
	\IEEEauthorblockA{\textit{Ingeniería en Computación} \\
		\textit{Facultad de Ingeniería}\\
		No. Cuenta: 314256009}
	\and
	\IEEEauthorblockN{Hernández Alonso Antonio Adonahí}
	\IEEEauthorblockA{\textit{Ingeniería en Computación} \\
		\textit{Facultad de Ingeniería}\\
		No. Cuenta: 320257780}
	\and
	\IEEEauthorblockN{Jiménez Reyes Alan Yael}
	\IEEEauthorblockA{\textit{Ingeniería en Computación} \\
		\textit{Facultad de Ingeniería}\\
		No. Cuenta: 320026726}
}

\maketitle

\begin{center}
	\vspace{-0.3cm}
	{\small
		\textbf{Práctica No.:} 4 \quad | \quad
		\textbf{Grupo:} 6 \quad | \quad
		\textbf{Fecha:} DD/MM/AAAA
	}
	\vspace{0.2cm}
\end{center}

\begin{abstract}
	Este reporte documenta el análisis del proceso de arranque (boot) de la Raspberry Pi 4, desde el bootloader almacenado en EEPROM hasta la inicialización completa del sistema operativo bajo el control de systemd. Se describe la identificación de las etapas del arranque, la inspección del bootloader, el análisis de logs con \texttt{journalctl}, la medición de tiempos de arranque con \texttt{systemd-analyze}, y la creación de servicios personalizados. El documento aborda técnicas de optimización del tiempo de arranque y su medición cuantitativa, demostrando cómo systemd gestiona la inicialización del sistema en Linux embebido.
\end{abstract}

\begin{IEEEkeywords}
	Boot, EEPROM, bootloader, systemd, journalctl, systemd-analyze, servicios, Raspberry Pi, init system
\end{IEEEkeywords}

\section{Objetivo}

\subsection{Objetivo General}
Comprender las etapas del proceso de arranque de la Raspberry Pi, desde el bootloader en EEPROM hasta el sistema completamente iniciado, utilizando las herramientas de diagnóstico de systemd para inspeccionar, medir y modificar el proceso de arranque.

\subsection{Objetivos Específicos}
\begin{itemize}
	\item Identificar las etapas del arranque de la Raspberry Pi, del bootloader al sistema completo.
	\item Verificar la versión del bootloader (EEPROM) sin necesidad de actualizarlo.
	\item Inspeccionar el log completo de arranque con \texttt{journalctl}.
	\item Medir y analizar el tiempo de arranque con \texttt{systemd-analyze}.
	\item Crear un servicio systemd propio que arranque automáticamente uno de los programas desarrollados en prácticas anteriores.
	\item Optimizar el tiempo de arranque, midiendo el efecto real de los cambios aplicados.
	\item Registrar los avances en el repositorio de Git del curso.
\end{itemize}

\section{Introducción}

El proceso de arranque de un sistema embebido es una secuencia compleja de etapas que transforman un hardware inerte en un sistema operativo funcional. En la Raspberry Pi 4, este proceso involucra múltiples componentes: el bootloader almacenado en EEPROM, el firmware de la GPU, el kernel de Linux, y el sistema de inicialización systemd.

\subsection{El Bootloader en EEPROM}

A diferencia de computadoras tradicionales que utilizan BIOS o UEFI almacenadas en la placa madre, la Raspberry Pi 4 almacena su bootloader en un chip EEPROM integrado en la placa. Este bootloader es lo primero que se ejecuta al aplicar energía al sistema y se encarga de localizar y cargar el siguiente componente del proceso de arranque desde la tarjeta microSD.

El bootloader en EEPROM puede actualizarse mediante el sistema operativo en ejecución, permitiendo correcciones de errores y nuevas características sin reemplazar hardware. La versión del bootloader determina capacidades como el arranque desde USB, configuraciones de red PXE, y la gestión de energía durante el arranque.

\subsection{Firmware de la GPU y Device Tree}

El segundo componente del arranque es el firmware de la GPU (VideoCore IV), almacenado en archivos como \texttt{start4.elf} en la partición boot de la microSD. Este firmware lee el archivo de configuración \texttt{config.txt}, inicializa el hardware básico (RAM, reloj del sistema, periféricos), y carga el kernel de Linux junto con el Device Tree.

El archivo \texttt{config.txt} permite configurar parámetros de hardware antes de que Linux arranque, como la habilitación de interfaces I2C y SPI (utilizada en la Práctica 3), overclocking, asignación de memoria a la GPU, y configuración de salidas de video. Esta etapa es fundamental porque establece el estado inicial del hardware que el kernel asumirá al tomar control.

\subsection{Kernel de Linux}

El kernel de Linux, usualmente almacenado como \texttt{kernel8.img} para arquitecturas de 64 bits, toma control del procesador ARM después de ser cargado por el firmware de la GPU. El kernel inicializa los subsistemas de gestión de memoria, scheduling de procesos, y drivers de dispositivos, monta el sistema de archivos raíz, y finalmente ejecuta el primer proceso de espacio de usuario: systemd (PID 1).

\subsection{systemd como Sistema de Inicialización}

systemd es el sistema de inicialización (init system) utilizado por la mayoría de distribuciones modernas de Linux, incluyendo Raspberry Pi OS. A diferencia de sistemas de init tradicionales basados en scripts secuenciales, systemd utiliza un modelo de dependencias paralelas que permite arrancar servicios simultáneamente cuando sus dependencias están satisfechas.

systemd gestiona servicios mediante archivos de unidad (\textit{unit files}) que definen cómo y cuándo deben iniciarse los servicios. Proporciona herramientas poderosas para el diagnóstico del arranque:

\begin{itemize}
	\item \texttt{journalctl} — Sistema de logging centralizado que captura todos los mensajes del kernel y servicios.
	\item \texttt{systemd-analyze} — Herramienta para medir y visualizar el tiempo de arranque de cada servicio.
	\item \texttt{systemctl} — Interfaz para gestionar servicios, habilitar/deshabilitar su arranque automático, y consultar su estado.
\end{itemize}

\subsection{Etapas del Arranque de la Raspberry Pi}

La tabla \ref{tab:boot_stages} resume las cuatro etapas principales del proceso de arranque:

\begin{table}[H]
	\centering
	\caption{Etapas del arranque de la Raspberry Pi.}
	\label{tab:boot_stages}
	\small
	\begin{tabular}{|c|p{5.5cm}|}
		\hline
		\textbf{\#} & \textbf{Etapa} \\
		\hline
		1 & Bootloader (EEPROM): busca firmware en microSD. \\
		\hline
		2 & Firmware GPU: lee \texttt{config.txt}, carga kernel. \\
		\hline
		3 & Kernel: inicializa SO, monta raíz, ejecuta systemd. \\
		\hline
		4 & systemd: arranca servicios según dependencias. \\
		\hline
	\end{tabular}
\end{table}

\subsection{Relación con Prácticas Anteriores}

En la Práctica 3, se modificó el archivo \texttt{/boot/firmware/config.txt} para habilitar las interfaces I2C y SPI mediante las líneas:

\begin{lstlisting}[language=bash]
dtparam=i2c_arm=on
dtparam=spi=on
\end{lstlisting}

Esta configuración no activa los buses ``mágicamente'' después de que Linux arranca, sino que modifica la Etapa 2 del arranque: el firmware de la GPU lee estas directivas y configura el hardware correspondiente antes de cargar el kernel. Linux hereda esta configuración y hace disponibles las interfaces \texttt{/dev/i2c-1} y \texttt{/dev/spidev0.0} durante su inicialización.

Esta práctica profundiza en el entendimiento de todo el proceso de arranque, permitiendo diagnosticar problemas, optimizar tiempos de inicio, y crear servicios personalizados que se ejecuten automáticamente al encender el sistema.

\section{Desarrollo}

\subsection{Material y Equipo}

Para la realización de esta práctica no se requieren componentes de hardware adicionales. Se utilizó únicamente:

\begin{itemize}
	\item Raspberry Pi 4 Model B con Raspberry Pi OS instalado (configurada en prácticas anteriores)
	\item Laptop con acceso SSH al sistema
	\item Toolchain de desarrollo en C (GCC) previamente instalado
	\item Repositorio Git del curso con prácticas anteriores
	\item Programa \texttt{sensor\_combinado.c} de la Práctica 3 (o alternativamente \texttt{leds.c} de la Práctica 2)
\end{itemize}

Esta práctica es enteramente de software y se enfoca en el análisis y optimización del proceso de arranque del sistema operativo.

\section{Resultados y Análisis}

Contenido pendiente: Presentar los resultados obtenidos de cada parte del desarrollo.

\section{Conclusiones}

\textbf{Belmont Muñoz Samuel}. \\
Conclusión pendiente.

\textbf{Domínguez Palacios Jesús Alejandro}. \\
Conclusión pendiente.

\textbf{Dorantes Mejía Emmanuel}. \\
Conclusión pendiente.

\textbf{García Olvera Emiliano}. \\
Conclusión pendiente.

\textbf{Hernández Alonso Antonio Adonahí}. \\
Conclusión pendiente.

\textbf{Jiménez Reyes Alan Yael}. \\
Conclusión pendiente.

\begin{thebibliography}{00}
	\bibitem{systemd_docs} Lennart Poettering et al. (2024). ``systemd System and Service Manager'', \textit{systemd Documentation}. Disponible en: https://www.freedesktop.org/wiki/Software/systemd/
	\bibitem{rpi_boot} Raspberry Pi Ltd. (2024). ``Raspberry Pi Documentation --- Boot Sequence''. Disponible en: https://www.raspberrypi.com/documentation/computers/raspberry-pi.html
	\bibitem{device_tree} The kernel development community. (2024). ``Device Tree Usage'', \textit{Linux Kernel Documentation}. Disponible en: https://www.kernel.org/doc/html/latest/devicetree/usage-model.html
	\bibitem{journalctl} Lennart Poettering. (2024). ``journalctl(1) - Linux man page''. Disponible en: https://man7.org/linux/man-pages/man1/journalctl.1.html
\end{thebibliography}

\end{document}
